% page 258 du fac-simile (image p-257 du scan)
% bloc 157.5 x 229.0 mm ; marge gauche 8.0 mm
\pgc{-12.700mm}{0.000mm}{
\l{\sou{jungar}. (trans.)\cel{2}Ligar (tir-bestii) a veturo, a plugilo ;}
\l{\cel{3}inteligar vagoni sur relvoyo. - L.}
\l{}
\l{\sou{junglo}. En India : plana lando, marshoza, kovrita ye kani}
\l{\cel{3}e ye bushi dika e dornoza. - EF.}
\l{}
\l{\sou{junio.}\cel{2}La monato sisesma di la yaro (segun la maniero mo-}
\l{\cel{3}derna.)\cel{2}- DEFIRS.}
\l{}
\l{\sou{juniora}. (La persono) qua naskis pos (frato). - EF.}
\l{}
\l{\sou{junipero}.\cel{2}(\sou{bot.)}\cel{1}Arbusto ek la familio \textquotedbl{}cipresinei\textquotedbl{}, di qua}
\l{\cel{3}la beri esas aromatoza. - L.\cel{1}\sou{juniperus}. - EFI.}
\l{}
\l{\sou{junko}. (\sou{bot.)}\cel{2}Genero tipa di la familio \textquotedbl{}junkacei\textquotedbl{}, planto}
\l{\cel{3}herbatra, rekta e flexebla, qua kreskas en la tereni humi-}
\l{\cel{3}da. -L.\cel{1}\sou{juncus}. - FIS.}
\l{}
\l{\sou{juntar}.\cel{2}(\sou{trans}.)I. Proximigar du kozi ye un ek olia extre-}
\l{\cel{3}maji rispektiva, por ke oli interkontaktez od intertenez.}
\l{\cel{3}- II. (\sou{aludante plura personi, trupi,e}\cel{1}c.)\cel{2}Irar unionar}
\l{\cel{3}(su) ad (altra personi,altra trupi). - EFIS.}
\l{}
\l{\sou{jupo}. Parto di la muliero-vesto, qua kovras de la zono til}
\l{\cel{3}la pedi. - FR.}
\l{}
\l{\sou{jurar}. (\sou{trans.}) Promisar, afirmar (ulo ad ulu) selektante (ge-}
\l{\cel{3}nerale) deo kom testo pri lo. - EFIS.}
\l{}
\l{\sou{juraso.}\cel{2}(\sou{geol.}) Tereno ye formaceso sekundara inter la stra-}
\l{\cel{3}to liasa e la tereno kretatra,quan onu renkontras preci-}
\l{\cel{3}pue en la Jura-monti (Francia). - DEF.}
\l{}
\l{\sou{jurio}. (en\cel{1}\sou{la tribunalo kriminal}a).Ensemblo ek la civitani}
\l{\cel{3}qui devas decidar ka la akuzato esas (o ne esas) kulpoza.}
\l{\cel{3}- DEFIS.}
\l{}
\l{\sou{juriskonsulto}. La persono qua konsilas profesione pri la}
\l{\cel{3}yuro-problemi. - DEFIRS.}
\l{}
\l{\sou{jurisprudenco.}\cel{2}Interpretado di la legi per la verdikti da}
\l{\cel{3}la judiciisti : o, generale, maniero aplikar la legi a}
\l{\cel{3}la kazi partikulara. - DFIS.}
\l{}
\l{\sou{jurnalo}.\cel{2}Edituro singladia o periodala, qua donas la\cel{2}niuzi}
\l{\cel{3}politikala, literaturala, ciencala. - DEFIS.}
\l{}
\l{\sou{jus}. Pasis nur kelka instanti de kande... - DEFIS.}
\l{}
\l{\sou{justa}.\cel{2}I.Qua precize fitas,su adaptas a to por quo olu des-}
\l{\cel{3}tinesis. (Dicesas pri mezuro,balanco, voco, vorto, penso).}
\l{\cel{3}- II. Qua agas o per quo onu agas perfekte lo exekutenda.}
\l{\cel{3}- III. Qua judikas korekte (en la kazo konsiderata). -EFIS.}
}
